% !TEX root = ../CerdaMS.tex

\section{Ohana trees}\label{sec:MemTrees}

% BEGIN unchanged author include/MTrees.tex
% !TEX root = ../CerdaMS.tex

\begin{figure}
	\centering
	\scalebox{.85}{%
	% MT of E_x
	\begin{tikzpicture}[baseline]
		\node[n] 	(E-lx1) 					{$\ensuremath{\lambda}x_1.x_1$};
		\node[left=0pt of E-lx1.west] {$\MT{\comb E_x} =$};
		\node[n] 	(E-bot1) 	at (-1,-1) 		{$\bot_\emptyset$};
		\node[n] 	(E-lx2) 	at (1.5,-1) 		{$\ensuremath{\lambda}x_2.x_2$};
		\node[n] 	(E-bot2) 	at (0.5,-2) 		{$\bot_{\set x}$};
		\node[n] 	(E-lx3) 	at (3,-2) 		{$\ensuremath{\lambda}x_3.x_3$};
		\node[n] 	(E-bot3) 	at (2,-3) 		{$\bot_\emptyset$};
		\node 		(E-inf) 	at (4.5,-3) 		{};
		\draw[b] ($(E-lx1.south)+(3pt,0)$) to[pos=.5] node[m]{$\emptyset$} (E-bot1);
		\draw[b] ($(E-lx1.south east)-(2pt,0)$) to[pos=.5] node[m]{$\set x$} (E-lx2);
		\draw[b] ($(E-lx2.south)+(3pt,0)$) to[pos=.5] node[m]{$\set x$} (E-bot2);
		\draw[b] ($(E-lx2.south east)-(2pt,0)$) to[pos=.5] node[m]{$\set x$} (E-lx3);
		\draw[b] ($(E-lx3.south)+(3pt,0)$) to[pos=.5] node[m]{$\emptyset$} (E-bot3);
		\draw[b,dotted] ($(E-lx3.south east)-(2pt,0)$) to[pos=.5] node[m]{$\set x$} (E-inf);
	\end{tikzpicture}
	% MT of O_x
	\quad
	\begin{tikzpicture}[baseline,overlay]
			\node[n] 	(O-lx1) 					{$\ensuremath{\lambda}x_1.x_1$};
		\node[left=0pt of O-lx1.west] {$\MT{\comb O_x} =$};
		\node[n] 	(O-bot1) 	at (-1,-1) 		{$\bot_{\set x}$};
		\node[n] 	(O-lx2) 	at (1.5,-1) 		{$\ensuremath{\lambda}x_2.x_2$};
		\node[n] 	(O-bot2) 	at (0.5,-2) 		{$\bot_\emptyset$};
		\node[n] 	(O-lx3) 	at (3,-2) 		{$\ensuremath{\lambda}x_3.x_3$};
		\node[n] 	(O-bot3) 	at (2,-3) 		{$\bot_{\set x}$};
		\node 		(O-inf) 	at (4.5,-3) 		{};
		\draw[b] ($(O-lx1.south)+(3pt,0)$) to[pos=.5] node[m]{$\set x$} (O-bot1);
		\draw[b] ($(O-lx1.south east)-(2pt,0)$) to[pos=.5] node[m]{$\set x$} (O-lx2);
		\draw[b] ($(O-lx2.south)+(3pt,0)$) to[pos=.5] node[m]{$\emptyset$} (O-bot2);
		\draw[b] ($(O-lx2.south east)-(2pt,0)$) to[pos=.5] node[m]{$\set x$} (O-lx3);
		\draw[b] ($(O-lx3.south)+(3pt,0)$) to[pos=.5] node[m]{$\set x$} (O-bot3);
		\draw[b,dotted] ($(O-lx3.south east)-(2pt,0)$) to[pos=.5] node[m]{$\set x$} (O-inf);
	\end{tikzpicture} 
	% MT of YRl
	\hspace{5cm}
	\begin{tikzpicture}[baseline]
		\node[n] 	(YR-l0) 					{$\ensuremath{\lambda}x_0.x_0$};
		\node[left=0pt of YR-l0.west] {\llap{$\MT{\comb Y\comb Rl} =$}};
		\node[n]	(YR-l1)		at (8pt,-1.1)		{$\ensuremath{\lambda}x_1.x_1$};
		\node[n]	(YR-l2)		at (16pt,-2.2)		{$\ensuremath{\lambda}x_2.x_2$};
		\node		(YR-inf)	at (24pt,-3.3)		{};
		\draw[b] ($(YR-l1)+(0,26pt)$) to[pos=.5] node[m]{$\set{l,x_0}$} (YR-l1);
		\draw[b] ($(YR-l2)+(0,26pt)$) to[pos=.5] node[m]{$\set{l,x_0,x_1}$} (YR-l2);
		\draw[b,dotted] ($(YR-inf)+(0,26pt)$) to[pos=.5] 
			node[m]{$\set{l,x_0,x_1,x_2}$} (YR-inf);
	\end{tikzpicture}%
	}
	\caption{Examples of Ohana trees of \lamI-terms.}\label{fig:MTrees}
\end{figure}
% END unchanged author include/MTrees.tex


We have seen that Böhm trees are not well-suited for representing \lamI-terms faithfully, because variables that are never erased along the reduction are forgotten (i.e., pushed into infinity).
We now introduce a notion of evaluation tree that records all variables that are pushed along each path, 
 ensuring that none are forgotten, whether they remain in passive position or not.

\subsection{The coinductive definition}

For every finite set $X\subseteq\Var$, written $X\subf \Var$, we introduce a constant $\bot_X$ representing any looping $M\in\LamI(X)$. 
Intuitively, the \emph{Ohana tree} of a \lamI-term $M$ is obtained from $\BT{M}$ by annotating each subtree (also 
$\bot$) with the free variables of the \lam-term that generated it.

\Needspace{3cm}
\begin{definition}\label{def:memorytrees4LamI}
The \emph{Ohana tree} of a term $M\in\LamI(X)$ is coinductively defined as follows:
\bsub
\item If $M \redd[h] \lam x_1\dots x_n.yM_1\cdots M_k$ with $y\in 
X\cup\set{\vec x}$, $M_i\in\LamI(X_i)$ (for $1\le i \le k$), then\smallskip
\[
\begin{tikzpicture}
\node 		(MTM) 	at (-22pt,0)				{$\MT M \eqdef$};
\node[n] 	(y) 	at (1.15,0) 		{$\lam x_1\dots x_n.y$};
\node[n] 	(MTM1) 	at (.6,-1) 	{$\MT{M_1}$};
\node 				at (2,-1) 		{\dots};
\node[n] 	(MTM2) 	at (3.4,-1) 	{$\MT{M_k}$};
\draw[b] ($(y.south east)+(-8pt,0)$) to[pos=0.5]node[m]{$X_1$} (MTM1.north);
\draw[b] (y.south east) to[pos=0.5]node[m]{$X_k$} (MTM2.north);
\end{tikzpicture}
\]
\item $\MT{M} \eqdef \bot_X$, otherwise. Recall that $X = \fv(M)$, 
since $M\in\LamI(X)$.
\esub
\end{definition}

We sometimes use an inline presentation of Ohana trees, 
by introducing application symbols annotated with the sets $X_i$, as in 
$\lam x_1\ldots x_n.y\cdot^{X_1}T_1\cdots^{X_k}T_k$.
\begin{example}\label{ex:MTs}\mbox{}
\bsub
\item $\MT{\comb{I}} = \lam x.x$, $\MT{\comb{D}} = \lam x.x\cdot^{\set{x}}x$, $\MT{\comb{B}} = \lam fgx.f\cdot^{\set{g,x}}(g\cdot^{\set{x}}x)$.
\item $\MT{\Om} = \MT{\lam y.\Om y} = \bot_{\emptyset}$, $\MT{\Om x} = \bot_{\set x}$. More generally: $\MT{\Om \vec x} = \bot_{\set{x_1,\dots,x_n}}$. 
\item  $\MT{\comb{Y}} = \lam f.f\cdot^{\set{f}}(f\cdot^{\set{f}}(f\cdot^{\set{f}}(\,\cdots)))$. 
	In fact, it is easy to check that for all \emph{closed} fpcs $Y\in\LamI$, we have $\MT{Y} = \MT{\comb{Y}} \cong \BT{\comb{Y}}$.
\item In $\MT{\Bible}$ the variable $l$ which is pushed into infinity becomes visible: $\MT{\Bible} =\MT{\Theta_l} = \lam f.f\cdot^{\set{f,l}}(f\cdot^{\set{f,l}}(f\cdot^{\set{f,l}}(\,\cdots)))$. Thus $\MT{\Bible} \neq \MT{\Theta_x}$ when $x\neq l$.
\esub
\end{example}

Other examples of Ohana trees are depicted in Figure~\ref{fig:MTrees}.
Note that the Ohana trees of $\comb{E}_x$ and $\comb{O}_x$ are now distinguished.
The interest of ${\comb{YR}l}$ is that it pushes into infinity an increasing amount of variables: 
\[\comb{YR}l =_\beta \lam x_0.x_0((\comb{YR})(x_0l)) =_\beta \lam x_0.x_0(\lam x_1.x_1((\comb{YR})(x_1(x_0l))))\, {=_\beta} \cdots\]
So---at the limit---infinitely many variables are pushed into infinity (but only $l$ is free).

\begin{remark} If $M\in\LamI$ has a $\beta$-nf $N$, then there is an isomorphism $\MT{M}\cong N$.
In fact, since $N$ is finite, it is possible to reconstruct all the labels in $\MT{M}$.  
%Formally, proceed by induction on the structure of $\MT{M}$ for $(\Rightarrow)$, and of $N$ for $(\Leftarrow)$.
\end{remark}

We are going to show that Ohana trees are invariant under $\beta$-reduction (Prop.~\ref{prop:InvariantUnderBeta} plus Thm.~\ref{thm:MTbijApp}), and that the equality induced on \lamI-terms is a \lamI-theory (Corollary~\ref{cor:MisALamITheory}).

\subsection{The associated approximation theory}

We introduce a theory of program approximation that captures Ohana trees, by adapting the well-established approach originally developed for Böhm trees~\cite{Levy76,Wadsworth78} (see also \cite[Ch.~14]{Bare}).
We start by defining the approximants with memory, corresponding to finite Ohana trees.

\begin{definition}\label{def:approximants}\mbox{}
\bsub
\item \label{def:approximants1}
	The set of \emph{approximants with memory} is 
	$\Appset \eqdef \bigdunion_{X\subf \Var} \Appset(X)$, 
	where $\Appset(X)$ is defined by (for $A_i\in\Appset(X_i)$, $X_i\subseteq X$, 
	$y\in X\cup\set{\vec x}$ and $\vec x\in\bigcup_{i=1}^k X_i$):
	\bnf[\Appset(X)]{A}{ \bot_X \mid \lam x_1\ldots x_n.yA_1\cdots A_k}
	
%\item We extend $\fv(\cdot)$ to $\Appset$ by adding $\fv(\bot_X) = X$, thus $\fv(A) = X$ stands for $A\in\Appset(X)$.
\item We define $\mathord{\MTle} \subseteq \Appset^2$ as 
the least partial order closed under the following rules:
\[
	\infer{\bot_X\MTle A}{A\in\Appset(X)}\qquad
	\infer{\lam x_1\dots x_n.y\, A_1\cdots A_k \MTle \lam x_1\dots x_n.y\, A'_1\cdots A'_k}{A_i\MTle A'_i&A_i,A'_i\in\Appset(X_i)&\textrm{for all }1\le i \le k&\vec x\in\set{y}\cup\big(\bigcup_{i=1}^k X_i\big)}
\]
\esub
\end{definition}

Each $(\Appset(X), \MTle)$ is a pointed poset (partially ordered set) with bottom element $\bot_X$, while the poset $(\Appset,\MTle)$ is not pointed because it has countably many minimal elements.
Note that we do not annotate the application symbols 
in $A\in\Appset$ with the sets of variables $X_i$, 
since the labels in its $\bot$-nodes carry enough information 
to reconstruct the finite Ohana tree associated with $A$.
%Contrarily to the case of memory tree, we do not annotate the application symbol in $A\in\LamI$, since the labels in its $\bot$-nodes carry enough information to reconstruct the finite memory tree associated with $A$.
%Remark that $(\Appset,\MTle)$ is not \emph{pointed} as a poset (partially ordered set), because it has countably many minimal elements $\bot_X$. 
%However, each $\Appset(X)$ has  its bottom element $\bot_X$.
%Also, we do not need to annotate the application symbol in $A\in\LamI$, 
%since the labels in its $\bot$-nodes carry enough information to 
%reconstruct 
%the finite memory tree associated with $A$.

\newcounter{lem:bijAppsetfinmemtrees}
\setcounter{lem:bijAppsetfinmemtrees}{\thetheorem}
\begin{lemma}
\label{lem:bijAppset_finmemtrees}\label{lem:bijAppsetfinmemtrees}
There is a bijection between $\Appset$ and the set of finite Ohana trees of \lamI-terms.
\end{lemma}

We now give the recipe to compute the set of approximants with memory of a \lamI-term.

\begin{definition}\mbox{}
\bsub
\item The \emph{direct approximant} of $M\in\LamI(X)$ is given by:
\[
	\dirapp{M} \eqdef 
	\begin{cases}
	 \lam x_1\ldots x_n.y\,\dirapp{M_1}\cdots \dirapp{M_k},&\textrm{if }M = \lam x_1\ldots x_n.yM_1\cdots M_k,\\	
	\bot_X,&\textrm{otherwise.}
	\end{cases}
\]
\item The set $\Appof{M}$ of \emph{approximants with memory of $M\in\LamI$} is defined by:
\[
	\Appof{M} \eqdef \set{ A\in\Appset \mid \exists N\in\LamI\,.\, M\redd[\beta]N \textrm{ and }A\MTle \dirapp{N}}
\]
\esub
\end{definition}

\begin{remark} \label{rem:appof-M-fv-M}
	For all $M \in \LamI(X)$, we have $\dirapp M \in \Appset(X)$ and $\Appof M \subseteq \Appset(X)$.
\end{remark}

\begin{example}\mbox{}
\bsub
\item $\Appof{\comb{D}} = \set{\bot_\emptyset, \lam x.x\bot_{\set x}, 
		\comb{D}}$, 
	$\Appof{\Om} = \set{\bot_{\emptyset}}$%
%	, $\Appof{\Om x} = \set{\bot_{\set{x}}}$%
	.
\item $\Appof{\comb{Y}} = \set{\bot_\emptyset} \cup 
		\set{ \lam f.f^n(\bot_{\set f}) \mid n\in\Nat} $.
\item $\Appof{\Bible} = \set{\bot_{\set{l}}} \cup 
		\set{ \lam f.f^n(\bot_{\set{f,l}}) \mid n\in\Nat} $.
\item $\Appof{\comb{YR}l} = \set{\bot_{\set{l}},
	\lam x_0.x_0(\bot_{\set{x_0,l}}),
	\lam x_0.x_0(\lam x_1.x_1(\bot_{\set{x_0,x_1,l}})), \dots}$.
\esub
\end{example}

\newcounter{lem:DirApp}
\setcounter{lem:DirApp}{\thetheorem}
\begin{lemma}\label{lem:DirApp}\mbox{}
\bsub
\item\label{lem:DirApp1} 
	For every $M,N\in\LamI$, $M\bred N$ entails  $\dirapp{M} \MTle \dirapp{N}$.
\item\label{lem:DirApp2} 
	For all $M\in\LamI(X)$, $\Appof{M}$ is an ideal of the poset $(\Appset(X),\sqsubseteq)$. More precisely:
	\begin{enumerate}
	\item
		non-emptiness: $\bot_X\in \Appof{M}$;
	\item
		downward closure: $A\in \Appof{M}$ and $A'\MTle A$ imply $A'\in\Appof{M}$;
	\item
		directedness: $\forall A_1,A_2\in\Appof{M},\exists A_3\in\Appof{M}$ such that $A_1\MTle A_3 \sqsupseteq A_2$.
	\end{enumerate}
\esub
\end{lemma}
\begin{proof}
\bsub
\item
By an easy induction on the structure of $M\in\LamI(X)$, using Remarks~\ref{rem:hnfs} and \ref{rem:aboutLambdaI}.
\item (a) and (b) are trivial. 
	For (c), define $\MTlub$ by setting: 
	$\bot_X\MTlub A = A\MTlub \bot_X= A$, $x\MTlub x = x$, 
	if $A_1,A'_1$ and $A_2, A'_2$ are compatible, 
	then $(\lam x.A_1)\MTlub (\lam x.A'_1) = \lam x.(A_1\MTlub A'_1)$ 
	and $A_1A_2\MTlub A'_1A'_2 = (A_1\MTlub A'_1) (A_2\MTlub A'_2)$. 
	Check that $A_3 = A_1 \MTlub A_2$ is their supremum.
\qedhere
\esub
\end{proof}

The next result shows that $\beta$-convertible terms share the same the set of approximants.\par\addvspace{1.5ex}

\noindent\begin{minipage}{\textwidth}
\newcounter{prop:InvariantUnderBeta}
\setcounter{prop:InvariantUnderBeta}{\thetheorem}
\begin{proposition}\label{prop:InvariantUnderBeta}
	Let $M,N\in\LamI(X)$. If $M\red[\beta] N$, then $\Appof{M} = \Appof{N}$.
\end{proposition}

\begin{proof}By confluence of $\red[\beta]$, using Lemma~\ref{lem:DirApp}\eqref{lem:DirApp1}.
\end{proof}
\end{minipage}

\newcounter{thm:MTbijApp}
\setcounter{thm:MTbijApp}{\thetheorem}
\begin{theorem}\label{thm:MTbijApp} 
%For all $M\in\LamI$, there is an isomorphism $\MT{M} \cong \Appof{M}$.
	There is a bijection $\mtapp{\cdot}$ between the set of Ohana trees and the set of approximants of \lamI-terms satisfying $\mtapp{\MT{M}} = \Appof{M}$, for all $M\in\LamI$.

\end{theorem}

\begin{proof} From $\MT{M}$ to $\Appof{M}$.
Let $\MT{M}^\star$ be the set of all finite subtrees of $\MT{M}$, 
where each truncated subtree $T$ is replaced by $\bot_{\fv(T)}$.
By Lemma~\ref{lem:bijAppset_finmemtrees}, $\MT{M}^\star\cong\Appof{M}$.

From $\Appof{M}$ to $\MT{M}$. For every finite path $\sigma$ in $\MT{M}$, there is an $A\in\Appof{M}$ representing $\MT{M}$ along $\sigma$. Conclude since $\MT{M}$ is the supremum of all such $\sigma$.
\end{proof}


